\chapter{迭代器}
\label{ch:17}

第 15 章你手写过一个 \texttt{filter}：自己声明 \texttt{List}、自己 \texttt{for-in}、自己 \texttt{push}。能跑，但每换一种操作——求平方、取前几个、求和——都得再手搓一遍循环。

这一章讲 \texttt{std:\allowbreak{}:\allowbreak{}iter}：它把这些套路做成了能随意拼的链。写起来像说话，读起来也像说话：

\par
\medskip
\begin{lstlisting}[style=chali]
std::iter::range(1, 100).filter(是偶数).map(乘十).take(3).collect()
\end{lstlisting}
\par
\medskip

\section{一条链由三段组成}

迭代器的写法永远是同一个形状：

\par
\medskip
\begin{lstlisting}[style=chali]
起点.中间操作.中间操作.….终结操作
\end{lstlisting}
\par
\medskip

\begin{itemize}

\item 起点：从哪儿开始——一个区间、一个列表；

\item 中间操作：\texttt{map} / \texttt{filter} / \texttt{take}……可以串好几个；

\item 终结操作：\texttt{collect} / \texttt{fold} / \texttt{count}……干活的是它。

\end{itemize}

要 import：

\par
\medskip
\begin{lstlisting}[style=chali]
import std::iter;
\end{lstlisting}
\par
\medskip

\section{起点：\texttt{range} 和 \texttt{List.\allowbreak{}iter(\allowbreak{})\allowbreak{}}}

最常用的起点是 \texttt{std:\allowbreak{}:\allowbreak{}iter:\allowbreak{}:\allowbreak{}range(\allowbreak{}start,\allowbreak{} stop)\allowbreak{}}——从 \texttt{start} 到 \texttt{stop}（不含 \texttt{stop}）：

\par
\medskip
\begin{lstlisting}[style=chali]
import std::iter;

void main()
{
    let xs = std::iter::range(1, 6).collect();
    println("{} {}", xs.len(), xs[4]);

}
\end{lstlisting}
\par
\medskip

\par
\medskip
\begin{lstlisting}[style=out]
5 5
\end{lstlisting}
\par
\medskip

\texttt{range(\allowbreak{}1,\allowbreak{} 6)\allowbreak{}} 数是 1、2、3、4、5——和切片一样左闭右开（第 12 章）。\texttt{collect(\allowbreak{})\allowbreak{}} 把它们收进一个 \texttt{List<\allowbreak{}int64>\allowbreak{}}。

还有一个重载可以指定步长：\texttt{range(\allowbreak{}start,\allowbreak{} stop,\allowbreak{} step)\allowbreak{}}。

另一种起点是从现成的容器出发：\texttt{列表.\allowbreak{}iter(\allowbreak{})\allowbreak{}}。它和 \texttt{range} 有一个重要区别——\texttt{range} 每次调用都造一条新链，而容器迭代器是一次性消费品（见 17.8）。

\section{中间操作：\texttt{map} / \texttt{filter} / \texttt{take} / \texttt{skip} / \texttt{rev}}

\texttt{map}——对每一项做同一个变换：

\par
\medskip
\begin{lstlisting}[style=chali]
import std::iter;

void main()
{
    let ys = std::iter::range(1, 6).map(fn(int64 x)
    {
        x * x
    }).collect();
    println("{} {} {}", ys[0], ys[1], ys[4]);

}
\end{lstlisting}
\par
\medskip

\par
\medskip
\begin{lstlisting}[style=out]
1 4 25
\end{lstlisting}
\par
\medskip

\texttt{filter}——只留下满足条件的：

\par
\medskip
\begin{lstlisting}[style=chali]
import std::iter;

void main()
{
    let evens = std::iter::range(1, 11)
    .filter(fn(int64 x)
    {
        x % 2 == 0
    })
    .collect();
    println("{} -> {} {} {}", evens.len(), evens[0], evens[1], evens[4]);

}
\end{lstlisting}
\par
\medskip

\par
\medskip
\begin{lstlisting}[style=out]
5 -> 2 4 10
\end{lstlisting}
\par
\medskip

\texttt{take(\allowbreak{}n)\allowbreak{}} / \texttt{skip(\allowbreak{}n)\allowbreak{}} 取前 n 个 / 跳过前 n 个；\texttt{rev(\allowbreak{})\allowbreak{}} 反转。它们和 \texttt{map} / \texttt{filter} 一样，都是「中间操作」，返回的还是迭代器，可以继续往下接。

\section{终结操作：\texttt{collect} / \texttt{fold} / \texttt{count}}

中间操作只是搭管道，真正让数据流动的是终结操作。常用的三个：

\begin{itemize}

\item \texttt{collect(\allowbreak{})\allowbreak{}} —— 收成一个 \texttt{List}；

\item \texttt{count(\allowbreak{})\allowbreak{}} —— 数一共多少个；

\item \texttt{fold(\allowbreak{}初值,\allowbreak{} 函数)\allowbreak{}} —— 从头到尾累积，最灵活。

\end{itemize}

\texttt{fold} 是「求和、求积、求最大」这类操作的通用形式。求和：

\par
\medskip
\begin{lstlisting}[style=chali]
import std::iter;

void main()
{
    int64 s = std::iter::range(1, 101).fold(0, fn(int64 acc, int64 x)
    {
        acc + x
    });
    println("{}", s);
    println("{}", std::iter::range(1, 101).count());

}
\end{lstlisting}
\par
\medskip

\par
\medskip
\begin{lstlisting}[style=out]
5050
100
\end{lstlisting}
\par
\medskip

\texttt{fold(\allowbreak{}0,\allowbreak{} …)\allowbreak{}} 里那个 \texttt{0} 是初值 \texttt{acc}，然后每来一个 \texttt{x}，用函数算出一个新的 \texttt{acc}：\texttt{0 → 0+1 → 1+2 → …}。想求最大值，初值给一个极小值，函数换成「取较大的」就行。

还有几个「问一句就走」的终结操作：

\par
\medskip
\begin{lstlisting}[style=chali]
import std::iter;

void main()
{
    println("{}", std::iter::range(1, 6).all(fn(int64 x)
    {
        x > 0
    }));
    println("{}", std::iter::range(1, 6).any(fn(int64 x)
    {
        x > 3
    }));
    println("{}", std::iter::range(1, 6).find(fn(int64 x)
    {
        x % 2 == 0
    }).unwrap());

}
\end{lstlisting}
\par
\medskip

\par
\medskip
\begin{lstlisting}[style=out]
1
1
2
\end{lstlisting}
\par
\medskip

\begin{itemize}

\item \texttt{all(\allowbreak{}pred)\allowbreak{}} 全满足吗；\texttt{any(\allowbreak{}pred)\allowbreak{}} 有满足的吗——都返回 \texttt{bool}；

\item \texttt{find(\allowbreak{}pred)\allowbreak{}} 返回第一个满足的（\texttt{Option}，第 16 章——所以后面能接 \texttt{.\allowbreak{}unwrap(\allowbreak{})\allowbreak{}}）。

\end{itemize}

\section{串成一条链}

现在把中间操作串起来看威力——「1 到 99 里，乘 2 之后能被 3 整除的前 3 个」：

\par
\medskip
\begin{lstlisting}[style=chali]
import std::iter;

void main()
{
    let ys = std::iter::range(1, 100)
    .map(fn(int64 x)
    {
        x * 2
    })
    .filter(fn(int64 x)
    {
        x % 3 == 0
    })
    .take(3)
    .collect();
    println("{} {} {} {}", ys.len(), ys[0], ys[1], ys[2]);

}
\end{lstlisting}
\par
\medskip

\par
\medskip
\begin{lstlisting}[style=out]
3 6 12 18
\end{lstlisting}
\par
\medskip

一行一句、从上往下读，和需求一一对应：\texttt{map} 乘 2、\texttt{filter} 要 3 的倍数、\texttt{take(\allowbreak{}3)\allowbreak{}} 取三个、\texttt{collect} 收起来。没有 \texttt{for-in}、没有下标、没有临时变量。

\section{惰性：写了 \ensuremath{\neq} 跑了}

这是迭代器最容易误解的地方。看下面这段：

\par
\medskip
\begin{lstlisting}[style=chali]
import std::iter;

int64 seen = 0;

void note(int64 x)
{
    seen = seen + 1;
}

void main()
{
    let lazy = std::iter::range(1, 6)
    .inspect(note)
    .map(fn(int64 x)
    {
        x * 10
    });
    println("{}", seen);

    let got = lazy.collect();
    println("{} {}", seen, got[0]);

}
\end{lstlisting}
\par
\medskip

\par
\medskip
\begin{lstlisting}[style=out]
0
5 10
\end{lstlisting}
\par
\medskip

\texttt{inspect} 是个「顺手看一眼」的中间操作——每过一项就调用一次 \texttt{note}。可注意：

\begin{itemize}

\item 第一次打印 \texttt{seen} 是 0——虽然链子搭好了，一项都还没走；

\item 直到 \texttt{.\allowbreak{}collect(\allowbreak{})\allowbreak{}} 那一刻，才真的从头跑到尾，\texttt{seen} 变成 5。

\end{itemize}

中间操作是惰性的：不接终结操作，它们一行都不执行。 这条设计有两个好处：

\begin{enumerate}

\item 省事——\texttt{take(\allowbreak{}3)\allowbreak{}} 时，源数据只要产出 3 个就够了，后面的根本不碰；

\item 能做无限序列——\texttt{std:\allowbreak{}:\allowbreak{}iter:\allowbreak{}:\allowbreak{}repeat(\allowbreak{}x)\allowbreak{}} 可以永远产下去，只要你用 \texttt{take} 限住。

\end{enumerate}

\begin{tip}{提示}
记住一句话：迭代器链只有碰到 \texttt{collect} / \texttt{fold} / \texttt{count} 这类终结操作才会开跑。
\end{tip}

\section{一个坑：\texttt{List.\allowbreak{}iter(\allowbreak{})\allowbreak{}} 的元素是引用}

\texttt{range} 产出的是实打实的 \texttt{int64}，所以 \texttt{fn(\allowbreak{}int64 x)\allowbreak{}} 就够了。但从列表出发时，\texttt{列表.\allowbreak{}iter(\allowbreak{})\allowbreak{}} 产出的是元素的只读引用，fn 的参数得写成引用：

\par
\medskip
\begin{lstlisting}[style=chali]
import std::iter;

void main()
{
    List<int64> xs;
    xs.push(10);
    xs.push(20);
    xs.push(30);
    int64 total = xs.iter().fold(0, fn(int64 acc, int64 const& x)
    {
        acc + x
    });
    println("{}", total);

}
\end{lstlisting}
\par
\medskip

\par
\medskip
\begin{lstlisting}[style=out]
60
\end{lstlisting}
\par
\medskip

写 \texttt{fn(\allowbreak{}int64 acc,\allowbreak{} int64 x)\allowbreak{}} 会报「类型实参推导冲突」——因为元素实际是\texttt{int64 const\&}。这是借用检查器在提醒你：\texttt{.\allowbreak{}iter(\allowbreak{})\allowbreak{}} 借用了那个列表，而不是把它复制一份（省内存，也符合第 7 章「引用」那一套）。

要原地改元素，用 \texttt{.\allowbreak{}iter\allowbreak{}Mut(\allowbreak{})\allowbreak{}}——它产出可写引用 \texttt{T mut\&}（第 7 章）：

\par
\medskip
\begin{lstlisting}[style=chali]
import std::iter;

void bump(int64 mut& x)
{
    x = x * 10;
}

void main()
{
    List<int64> xs;
    xs.push(1);
    xs.push(2);
    xs.push(3);
    xs.iterMut().forEach(bump);
    println("{} {} {}", xs[0], xs[1], xs[2]);

}
\end{lstlisting}
\par
\medskip

\par
\medskip
\begin{lstlisting}[style=out]
10 20 30
\end{lstlisting}
\par
\medskip

\texttt{bump} 拿到的是可写引用，所以它改的就是列表里的原件。

不想借、想要一份独立的元素值，用 \texttt{.\allowbreak{}into\allowbreak{}Iter(\allowbreak{})\allowbreak{}}——它拿走列表的元素，fn 参数写 \texttt{int64} 就行。

\begin{center}
\begin{tabular}{lll}
\toprule
你想要 & 用 & 元素类型 \\
\midrule
只读遍历 & \texttt{xs.\allowbreak{}iter(\allowbreak{})\allowbreak{}} & \texttt{T const\&} \\
边走边改元素 & \texttt{xs.\allowbreak{}iter\allowbreak{}Mut(\allowbreak{})\allowbreak{}} & \texttt{T mut\&} \\
拿走元素 & \texttt{xs.\allowbreak{}into\allowbreak{}Iter(\allowbreak{})\allowbreak{}} & \texttt{T} \\
\bottomrule
\end{tabular}
\end{center}

\section{迭代器是一次性消费品}

17.2 留的那个问题现在回答：\texttt{range} 和 \texttt{列表.\allowbreak{}iter(\allowbreak{})\allowbreak{}} 的区别是——迭代器只能从头到尾走一遍。走完之后再调 \texttt{next(\allowbreak{})\allowbreak{}}、再接 \texttt{all} / \texttt{any} / \texttt{fold}，得到的都是空：不会报错，就是没有数据。

\par
\medskip
\begin{lstlisting}[style=chali]
import std::iter;

void main()
{
    List<int64> xs;
    xs.push(1);
    xs.push(2);
    xs.push(3);
    xs.push(4);
    let first = xs.iter();
    println("第一次 all: {}", first.all(fn(int64 const& x)
    {
        x > 0
    }));
    println("第二次 any: {}", first.any(fn(int64 const& x)
    {
        x > 2
    }));
    let second = xs.iter();
    println("新建后 any: {}", second.any(fn(int64 const& x)
    {
        x > 2
    }));

}
\end{lstlisting}
\par
\medskip

\par
\medskip
\begin{lstlisting}[style=out]
第一次 all: 1
第二次 any: 0
新建后 any: 1
\end{lstlisting}
\par
\medskip

\texttt{first} 被 \texttt{all} 走完就耗干了，第二次 \texttt{any} 是在空迭代器上问的——答案当然是「没有」。想再走一遍，就再调一次 \texttt{xs.\allowbreak{}iter(\allowbreak{})\allowbreak{}} 造一条新的（\texttt{second} 那两行）。\texttt{range} 没这个烦恼，因为它每次调用本来就是一条新链。

记住一句话：迭代器是一次性的，链子要走第二遍就重新造起点。

\subsection*{从列表造起点：\texttt{from\allowbreak{}List}}

\texttt{std:\allowbreak{}:\allowbreak{}iter:\allowbreak{}:\allowbreak{}from\allowbreak{}List<\allowbreak{}T>\allowbreak{}(\allowbreak{}list)\allowbreak{}} 也是一条起点：把列表的元素复制进迭代器，原来的列表完好无损：

\par
\medskip
\begin{lstlisting}[style=chali]
import std::iter;

void main()
{
    List<int64> xs;
    xs.push(2);
    xs.push(4);
    xs.push(6);
    int64 total = fromList<int64>(xs).fold(0, fn(int64 acc, int64 x)
    {
        acc + x
    });
    println("{} {}", total, xs.len());

}
\end{lstlisting}
\par
\medskip

\par
\medskip
\begin{lstlisting}[style=out]
12 3
\end{lstlisting}
\par
\medskip

\subsection*{\texttt{Set.\allowbreak{}iter(\allowbreak{})\allowbreak{}} 与 \texttt{Map.\allowbreak{}iter(\allowbreak{})\allowbreak{}}}

\texttt{Set<\allowbreak{}T>\allowbreak{}} 和 \texttt{Map<\allowbreak{}K,\allowbreak{}V>\allowbreak{}} 同样有 \texttt{.\allowbreak{}iter(\allowbreak{})\allowbreak{}}：借用快照——把元素（Map 只能是 \texttt{Pair<\allowbreak{}K,\allowbreak{} V>\allowbreak{}}，用 \texttt{first} 读取键、\texttt{second} 读取值）复制进迭代器，来源容器保持不变；对应的 \texttt{.\allowbreak{}into\allowbreak{}Iter(\allowbreak{})\allowbreak{}} 则把容器搬空。

\par
\medskip
\begin{lstlisting}[style=chali]
import std::iter;

void main()
{
    Set<int64> s;
    s.add(1);
    s.add(2);
    s.add(3);
    let si = s.iter();
    int64 sum = 0;
    loop
    {
        match(si.next())
        {
            Value(v) => { sum = sum + v; }
            None => { break; }
        }
    }
    println("sum = {} len = {}", sum, s.len());
    Map<int64, int64> m;
    m.put(1, 10);
    m.put(2, 20);
    let mi = m.iter();
    int64 doubled = 0;
    loop
    {
        match(mi.next())
        {
            Value(pair) => { doubled = doubled + pair.first * pair.second; }
            None => { break; }
        }
    }
    println("doubled = {} len = {}", doubled, m.len());

}
\end{lstlisting}
\par
\medskip

\par
\medskip
\begin{lstlisting}[style=out]
sum = 6 len = 3
doubled = 50 len = 2
\end{lstlisting}
\par
\medskip

迭代完了 \texttt{s.\allowbreak{}len(\allowbreak{})\allowbreak{}} 还是 3、\texttt{m.\allowbreak{}len(\allowbreak{})\allowbreak{}} 还是 2——借用快照不动原件。Map 的每一项是元组，用 \texttt{pair.\allowbreak{}0} 取键、\texttt{pair.\allowbreak{}1} 取值。

\section{小结}

{\small\setlength{\tabcolsep}{4pt}
\begin{longtable}{>{\raggedright\arraybackslash}p{\dimexpr 0.081\linewidth-2\tabcolsep\relax}>{\raggedright\arraybackslash}p{\dimexpr 0.899\linewidth-2\tabcolsep\relax}}
\toprule
你会了 & 关键点 \\
\midrule
\endhead
三段式 & 起点 \texttt{.\allowbreak{}} 中间操作 \texttt{.\allowbreak{}} … \texttt{.\allowbreak{}} 终结操作 \\
起点 & \texttt{std:\allowbreak{}:\allowbreak{}iter:\allowbreak{}:\allowbreak{}range(\allowbreak{}a,\allowbreak{} b)\allowbreak{}}（左闭右开）、\texttt{xs.\allowbreak{}iter(\allowbreak{})\allowbreak{}}、\texttt{from\allowbreak{}List(\allowbreak{}xs)\allowbreak{}}、\texttt{set.\allowbreak{}iter(\allowbreak{})\allowbreak{}} / \texttt{map.\allowbreak{}iter(\allowbreak{})\allowbreak{}}、\texttt{xs.\allowbreak{}into\allowbreak{}Iter(\allowbreak{})\allowbreak{}} \\
中间操作 & \texttt{map} \texttt{filter} \texttt{take} \texttt{skip} \texttt{rev} \texttt{inspect} \texttt{enumerate} \texttt{chain} \texttt{zip} \\
终结操作 & \texttt{collect} \texttt{fold} \texttt{count} \texttt{all} \texttt{any} \texttt{find} \texttt{last} \\
组合 & 中间操作可无限串联，从上往下读即需求 \\
惰性 & 不接终结操作就不执行；\texttt{take} 能省掉多余计算 \\
一次性 & 迭代器只能走一遍；要走第二遍就重新造起点 \\
引用坑 & \texttt{.\allowbreak{}iter(\allowbreak{})\allowbreak{}} 元素是 \texttt{T const\&}；要独立值用 \texttt{.\allowbreak{}into\allowbreak{}Iter(\allowbreak{})\allowbreak{}} \\
引入 & \texttt{import std:\allowbreak{}:\allowbreak{}iter;} \\
\bottomrule
\end{longtable}
}

\section{下一步}

单文件程序写多了，总得拆。函数、结构体、常量放哪儿？怎么让别的文件用上？测试和格式化怎么跑？

第 18 章讲模块与工程：\texttt{import} 到底导入了什么、\texttt{chali new} 起出来的工程长什么样、\texttt{chali test} / \texttt{chali fmt} 怎么用。

翻到\hyperref[ch:18]{第 18 章 模块与工程}。

\section{边界与验证}

惰性迭代器通常在消费时执行工作。修改当前元素与改变集合结构不是同一操作；遍历中的引用、源集合存活和结构修改按具体接口规则处理。

\section{习题}

前两题用于理解和基本实现，第三题检查边界，最后一题为综合挑战。先独立完成，再核对参考解答。

\begin{enumerate}

\item 构造一个惰性 map 链但不消费，能否断言映射工作已完成？

\item 使用 range、filter 和 collect 得到指定区间中的偶数。

\item 比较修改遍历元素与向源容器追加元素的区别。

\item 挑战：过滤集合后删除原集合中的目标项，设计不依赖失效索引的方法。

\end{enumerate}

\subsection*{参考解答与验收要点}

\begin{enumerate}

\item 不能；先查接口的惰性约定，再通过消费操作观察实际结果。

\item 明确区间端点，验证收集顺序和数量。

\item 前者可能通过元素可写引用发生，后者改变结构，不能据此前者合法就推出后者也合法。

\item 先收集稳定标识或构造新集合，再按明确策略更新；验证相邻目标、空集合和全匹配。

\end{enumerate}

\section{番外：少了一页的遍历}

改编自《西游记》通天河晒经，以下为编程主题的独立改写。

师徒过河，经书落水，沙僧把湿页依次摊开，八戒负责一边清点一边收走晒干的书。他刚数完第一本就将它拿走，下一本往前挪，他却照旧去看下一个位置。

到傍晚，清单上少了几本。八戒说每一步都严格递增，绝不会退步。悟空把书重新排好，让他观察：位置递增的同时，前面的集合也在缩短，两次变化合起来，恰好跳过了等待清点的经书。

他们改成先记录要收的书，再统一搬走。遇到接口不允许结构修改的地方，就先完成遍历。沙僧将逐项处理的次序写在纸上，不再只盯着最后那个计数。

唐僧问能否用更短的一条链完成。悟空说可以研究，但惰性链只有真正被消费才会工作。若只是把纸条写得漂亮，湿经书仍躺在河边，半页也不会自己翻干。

\subsection*{技术对照}

惰性迭代器通常在消费时执行工作。修改当前元素与改变集合结构不是同一操作；遍历中的引用、源集合存活和结构修改按具体接口规则处理。
